% problem-000310 11 石竹烯结构推断
\section*{11. 石竹烯结构推断}
反应1：

$
\text { Caryophyllene } \xrightarrow {\mathrm{H} _ {2} , \mathrm{Pd/C}} \mathrm{C} _ {1 5} \mathrm{H} _ {2 8}
$

反应 2：

（图：）

反应 3：

$
\text { Caryophyllene } \xrightarrow {\text { 1   eq   BH } _ {3} , \text { THF }} \xrightarrow {\mathrm{H} _ {2} \mathrm{O} _ {2} , \mathrm{NaOH} , \mathrm{H} _ {2} \mathrm{O}} \quad \mathbf {C} (\mathrm{C} _ {1 5} \mathrm{H} _ {2 6} \mathrm{O})
$

反应 4：

（图：）
D

⽯⽵烯异构体——异⽯⽵烯在反应1和反应2中也分别得到产物A和 $\mathbf { B } ,$ ，⽽在经过反应3后却得到了产物C的异构体，此异构体在经过反应4后仍得到了产物D。

\subsection*{11-1}
在不考虑反应⽣成⼿性中⼼的前提下，画出化合物A、C以及C的异构体的结构式。

\subsection*{11-2}
画出⽯⽵烯和异⽯⽵烯的结构式。

\subsection*{11-3}
指出⽯⽵烯和异⽯⽵烯的结构差别。
